\subsection{稳定丛的存在性}\label{src:9.2.5-subsection}
我们证明稳定丛的存在性 (\bookref{9.2.16}), 推广先前关于半稳定丛存在性的结果 (\bookref{9.1.36}). 论证采用模理论方法, 依赖于 $\cBun_{r,d}^{\mathrm s}$, $\cBun_{r,d}^{\mathrm{ss}}$ 及扩张叠的代数性.
\begin{theorem}[稳定丛的存在性]\label{src:9.2.16}
设 $C$ 是代数闭域 $\Bbbk$ 上亏格 $g\geq2$ 的光滑连通射影曲线. 则 $\cBun_{r,d}^{\mathrm s}$ 非空, 且在 $\cBun_{r,d}^{\mathrm{ss}}$ 中稠密.
\end{theorem}
\begin{proof}
定理~\bookref{9.1.36} 给出半稳定向量丛的存在性, 当 $\gcd(r,d)=1$ 时也给出稳定向量丛的存在性. 对秩作归纳, 通过证明 $\dim(\cBun_{r,d}^{\mathrm{ss}}\setminus\cBun_{r,d}^{\mathrm s})\leq\dim\cBun_{r,d}^{\mathrm{ss}}$ 来证明定理. 若 $\gcd(r,d)>1$, 则秩为 $r$, 次数为 $d$ 的向量丛 $F$ 严格半稳定, 当且仅当它可写为半稳定向量丛的扩张 $0\to F_2\to F\to F_1\to0$, 其中 $\operatorname{rk}(F_i)=r_i<r$, $\mu(F_i)=\mu(F)$. 令 $\underline{\mathrm{EXT}}(p_{12}^*\mathcal F_{1,\mathrm{univ}},p_{13}^*\mathcal F_{2,\mathrm{univ}})$ 为 $\cBun_{r_1,d_1}^{\mathrm{ss}}\times\cBun_{r_2,d_2}^{\mathrm{ss}}$ 上的扩张叠 (命题~\bookref{9.2.13}), 万有扩张 $\mathcal E_{\mathrm{univ}}$ (注记~\bookref{9.2.15}) 定义可表态射
\[
\underline{\mathrm{EXT}}(p_{12}^*\mathcal F_{1,\mathrm{univ}},p_{13}^*\mathcal F_{2,\mathrm{univ}})\to\cBun_{r,d}^{\mathrm{ss}},
\]
其像由包含秩为 $r_2$, 斜率为 $d/r$ 的子丛的严格半稳定向量丛 $E$ 组成. 定理由如下计算得到:
\[
\begin{aligned}
&\dim\underline{\mathrm{EXT}}(p_{12}^*\mathcal F_{1,\mathrm{univ}},p_{13}^*\mathcal F_{2,\mathrm{univ}})\\
&\quad=\dim\cBun_{r_1,d_1}^{\mathrm{ss}}+\dim\cBun_{r_2,d_2}^{\mathrm{ss}}\\
&\qquad\quad+\dim\operatorname{Ext}_{\mathcal O_C}^1(F_1,F_2)-\dim\operatorname{Ext}_{\mathcal O_C}^0(F_1,F_2)\\
&\quad=(r_1^2+r_2^2+r_1r_2)(g-1)+(d_1r_2-d_2r_1)\\
&\quad=(r_1^2+r_2^2+r_1r_2)(g-1)\\
&\quad<(r_1+r_2)^2(g-1)=\dim\cBun_{r,d}^{\mathrm{ss}}.
\end{aligned}
\]
第一步对任意 $F_i\in\cBun_{r_i,d_i}^{\mathrm{ss}}$ 成立, 第二步由 Riemann--Roch 定理 (\bookref{9.1.4}) 得到, 第三步使用 $d_1/r_1=d_2/r_2$.
\end{proof}
\begin{exercise}[简单]\label{src:9.2.17}
给出半稳定向量丛存在性 (\bookref{9.2.16}) 的另一种论证.

\textit{提示:} 遵循上述证明, 利用扩张叠以及不稳定化子丛的秩与次数只有有限种可能这一观察. 也可以不使用扩张叠, 而像 \cite[引理 4.3]{NR69} 与 \cite[\S8.6]{LP97} 那样利用 Quot 概形的维数来论证.
\end{exercise}
\begin{exercise}[简单]\label{src:9.2.18}
设 $r>1$ 和 $d$ 为整数.
\begin{enumerate}[(a)]
\item 证明 $\operatorname{codim}(\cBun_{r,d}\setminus\cBun_{r,d}^{\mathrm{ss}})\geq2$ 且 $\operatorname{codim}(\cBun_{r,\xi}\setminus\cBun_{r,\xi}^{\mathrm{ss}})\geq2$.
\item 证明若 $r\ne2$ 或 $g\ne2$, 则 $\operatorname{codim}_{\cBun_{r,d}^{\mathrm{ss}}}(\cBun_{r,d}^{\mathrm{ss}}\setminus\cBun_{r,d}^{\mathrm s})\geq2$; 若 $r=g=2$, 则余维为 $1$. 对固定行列式的情形证明类似结果.
\end{enumerate}
\end{exercise}
\subsection{单向量丛及其刚性化}\label{src:9.2.6-subsection}
\begin{definition}[单向量丛]\label{src:9.2.19}
$C$ 上向量丛 $E$ 称为\newterm{单}{simple}, 若 $\operatorname{Aut}(C)\cong\Bbbk^*$. 对 $\Bbbk$-概形 $S$, $C_S$ 上以 $S$ 为基底的向量丛族 $\mathcal E$ 称为单的, 若对每个 $s\in S$, 几何纤维 $\mathcal E_{\overline s}:=\mathcal E|_{C\times\Spec\kappa(\overline s)}$ 是单的.
\end{definition}
由引理~\bookref{9.1.15}, 每个稳定向量丛都是单的.
\begin{proposition}\label{src:9.2.20}
设 $\mathcal E$ 是 $\Bbbk$-概形 $S$ 上 $C_S$ 上的向量丛族.
\begin{enumerate}[(1)]
\item 存在群概形的闭浸入 $i_{\mathcal E}:\mathbb G_{m,S}\hookrightarrow\underline{\operatorname{Aut}}(\mathcal E)$, 与基变换相容: 若 $S'\to S$ 是概形映射, $\mathcal E'$ 表示 $\mathcal E$ 在 $C_{S'}$ 上的拉回, 则在自然同构 $\underline{\operatorname{Aut}}(\mathcal E)\times_SS'\cong\underline{\operatorname{Aut}}(\mathcal E')$ 下, $i_{\mathcal E'}$ 与 $i_{\mathcal E}$ 沿 $S'\to S$ 的基变换等同.
\item 若 $\mathcal E$ 是单的, 则 $i_{\mathcal E}:\mathbb G_{m,S}\xrightarrow{\sim}\underline{\operatorname{Aut}}(\mathcal E)$ 是同构.
\end{enumerate}
\end{proposition}
\begin{proof}
函子性地定义 $i_{\mathcal E}$: 对 $S$-概形 $S'$, 令 $i_{\mathcal E}(S'):\mathbb G_{m,S}(S')\to\operatorname{Aut}(\mathcal E')$ 为将单位元 $f\in\Gamma(S',\mathcal O_{S'})^*$ 映到由乘以单位元 $p_2^*f\in\Gamma(C_{S'},\mathcal O_{C_{S'}})^*$ 给出的同构 $\mathcal E'\to\mathcal E'$ 的映射. 因为 $i_{\mathcal E}(S')$ 单射, $i_{\mathcal E}:\mathbb G_{m,S}\to\operatorname{Aut}(\mathcal E)$ 是群概形的单态射. 为说明 $i_{\mathcal E}$ 是闭浸入, 先断言代数空间商 $X:=\underline{\operatorname{Aut}}(\mathcal E)/\mathbb G_{m,S}$ 在 $S$ 上仿射. 事实上, 因为 $p:\underline{\operatorname{Aut}}(\mathcal E)\to S$ 仿射且 $\mathbb G_m$ 线性约化, 自然映射 $X\to\Spec_S(p_*\mathcal O_X)^{\mathbb G_{m,S}}$ 是良模空间. 由定理~\bookref{7.5.11}, 此映射对到代数空间的映射万有, 因而是同构. 包含映射 $\mathbb G_{m,S}\to\underline{\operatorname{Aut}}(\mathcal E)$ 由 $X\to X\times X$ 沿 $\underline{\operatorname{Aut}}(\mathcal E)\to X\times X$ 作基变换得到. 因为 $X$ 分离, $\mathbb G_{m,S}\to\underline{\operatorname{Aut}}(\mathcal E)$ 是闭浸入.

对 (2), 首先注意纤维平坦性判据 (\bookref{A.2.10}) 蕴含 $i_{\mathcal E}$ 平坦. 因为 $i_{\mathcal E}$ 是有限型平坦满射单态射, 它是同构 (命题~\bookref{A.2.15}).
\end{proof}
\begin{exercise}\label{src:9.2.21}
对 $C$ 上向量丛 $E$, 证明下列条件等价:
\[
E\text{ 是单的}\iff\underline{\operatorname{Aut}}(E)\cong\mathbb G_m
\iff\dim\underline{\operatorname{Aut}}(E)=1.
\]
\end{exercise}
\begin{definition}\label{src:9.2.22}
秩为 $r$, 次数为 $d$ 的单向量丛叠为子叠
\[
\cBun_{r,d}^{\mathrm{simple}}\subseteq\cBun_{r,d},
\]
它在 $\Bbbk$-概形 $S$ 上的对象为秩 $r$, 次数 $d$ 的单向量丛族.
\end{definition}
\begin{proposition}\label{src:9.2.23}
子叠 $\cBun_{r,d}^{\mathrm{simple}}\subseteq\cBun_{r,d}$ 开, 且是代数空间 $M_{r,d}^{\mathrm{simple}}:=\cBun_{r,d}^{\mathrm{simple}}\!\sslash\mathbb G_m$ 上带 $\mathbb G_m$-带的 gerbe.
\end{proposition}
\begin{proof}
对开放性, 需证明若 $\mathcal E$ 是 $\Bbbk$-概形 $S$ 上的向量丛族, 则轨迹
\[
\{s\in S\mid\text{几何纤维 }\mathcal E_{\overline s}\text{ 是单的}\}
\]
是 $S$ 的开子集. 事实上, 由习题~\bookref{9.2.21}, 此轨迹是使群概形 $\operatorname{Aut}(\mathcal E)\to S$ 的纤维维数为 $1$ 的点 $s\in S$ 的集合; 由群概形的纤维维数的上半连续性 (命题~\bookref{B.1.9}), 它是开的. 命题~\bookref{9.2.20}(2) 的典范同构 $i_{\mathcal E}:\mathbb G_{m,S}\xrightarrow{\sim}\underline{\operatorname{Aut}}(\mathcal E)$ 诱导 $\cBun_{r,d}^{\mathrm{simple}}$ 的惯性叠与 $\mathbb G_m\times\cBun_{r,d}^{\mathrm{simple}}$ 之间的同构. 因而由命题~\bookref{7.4.19} (或命题~\bookref{7.4.55}) 对 $\mathbb G_m$ 作刚性化, 得到结论.
\end{proof}
\begin{remark}[向量丛叠的刚性化]\label{src:9.2.24}
可以应用命题~\bookref{7.4.55}, 利用惯性叠 $I_{\cBun_{r,d}}$ 的 fppf 闭子群概形 $\mathbb G_m:=\mathbb G_{m,\cBun_{r,d}}$, 构造所有向量丛叠的刚性化 $\cBun_{r,d}\!\sslash\mathbb G_m$. 该子群由命题~\bookref{9.2.20}(1) 的闭浸入 $i_{\mathcal E}:\mathbb G_{m,S}\hookrightarrow\underline{\operatorname{Aut}}(\mathcal E)$ 定义. 另一种方法是将 $\cBun_{r,d}=\bigcup_N[Q_N'/\mathrm{GL}_{P(N)}]$ 看作商叠之并 (定理~\bookref{4.1.23}), 其中 $Q_N'\subseteq\Quot^P(\mathcal O_C(-N)^{\oplus P(N)}/C)$ 是开子概形; 此时 $\cBun_{r,d}\!\sslash\mathbb G_m=\bigcup_N[Q_N'/\mathrm{PGL}_{P(N)}]$. 叠 $\cBun_{r,d}\!\sslash\mathbb G_m$ 以代数空间 $M_{r,d}^{\mathrm{simple}}$ 为开子叠.
\end{remark}
\begin{exercise}\label{src:9.2.25}
证明 $\cBun_{r,d}^{\mathrm{simple}}$ 不有界.
\end{exercise}
\textbf{应用于稳定丛.} 因为 $\cBun_{r,d}^{\mathrm s}\subseteq\cBun_{r,d}^{\mathrm{simple}}$ 是开子叠, 可应用命题~\bookref{9.2.23} 得到:
\begin{corollary}\label{src:9.2.26}
稳定丛叠 $\cBun_{r,d}^{\mathrm s}$ 是代数空间 $M_{r,d}^{\mathrm s}:=\cBun_{r,d}^{\mathrm s}\!\sslash\mathbb G_m$ 上带 $\mathbb G_m$-带的 gerbe.
\end{corollary}
\subsection{固定行列式的向量丛}\label{src:9.2.7-subsection}
\begin{definition}\label{src:9.2.27}
对线丛 $\xi\in\mathrm{Pic}(C)$, 定义 $(\Sch/\Bbbk)_{\text{ét}}$ 上的叠 $\cBun_{r,\xi}$, 其在 $\Bbbk$-概形 $S$ 上的对象为对 $(E,\varphi)$, 其中 $E$ 是 $C_S$ 上的向量丛, $\varphi:\det E\xrightarrow{\sim}p_1^*\xi$ 是同构. 态射 $(E,\varphi)\to(E',\varphi')$ 为与同构 $\varphi,\varphi'$ 相容的同构 $\alpha:E\xrightarrow{\sim}E'$.
\end{definition}
行列式为 $\xi$ 的向量丛叠 $\cBun_{r,\xi}$ 嵌入 Cartesian 图
\begin{equation}\label{src:9.2.28}
\begin{tikzcd}[ampersand replacement=\&]
\cBun_{r,\xi}\arrow[r]\arrow[d]\&\Spec\Bbbk\arrow[d,"\xi"]\\
\cBun_{r,d}\arrow[r,"\det"']\&\underline{\mathcal P\!ic}^{d},
\end{tikzcd}
\end{equation}
其中 $\underline{\mathcal P\!ic}^d=\cBun_{1,d}$ 是次数为 $d$ 的线丛的 Picard 叠 (另见 \S\bookref{7.4.8-subsection} 的讨论), $\det:\cBun_{r,d}\to\underline{\mathcal P\!ic}^d$ 是由 $E\mapsto\det E$ 定义的态射. 有开子叠
\[
\cBun_{r,\xi}^{\mathrm s}\subseteq\cBun_{r,\xi}^{\mathrm{ss}}\subseteq\cBun_{r,\xi},
\]
由行列式为 $\xi$ 的稳定及半稳定向量丛组成. 对秩为 $r$ 的稳定丛 $E$, 有典范同构 $\underline{\operatorname{Aut}}(E)\cong\boldsymbol\mu_r$. 特别地, $\cBun_{r,\xi}^{\mathrm s}$ 是 Deligne--Mumford 叠. 此外, 与命题~\bookref{9.2.23} 相同的论证说明 $\cBun_{r,\xi}^{\mathrm s}$ 是 $M_{r,\xi}^{\mathrm s}:=\cBun_{r,\xi}^{\mathrm s}\!\sslash\boldsymbol\mu_r$ 上的 $\boldsymbol\mu_r$-gerbe.
\begin{proposition}\label{src:9.2.29}
对每个线丛 $\xi\in\mathrm{Pic}(C)$ 和秩 $r>0$, 都存在 $C$ 上秩为 $r$, 满足 $\det(E)\cong\xi$ 的稳定向量丛 $E$.
\end{proposition}
\begin{proof}
由稳定丛的存在性 (\bookref{9.2.16}), 存在秩 $r$, 次数 $d=\deg\xi$ 的稳定丛 $E$. 则 $\xi\otimes(\det E)^{-1}$ 是次数 $0$ 线丛. 因为 $\operatorname{Pic}^0(C)$ 是 Abel 簇, 它是可除群, 因而存在次数 $0$ 线丛 $L$, 满足 $L^{\otimes r}\cong\xi\otimes(\det E)^{-1}$. 向量丛 $E\otimes L$ 半稳定 (引理~\bookref{9.1.13}(4)), 秩为 $r$, 且 $\det(E\otimes L)\cong(\det E)\otimes L^{\otimes r}\cong\xi$.
\end{proof}
\begin{exercise}[有益练习]\label{src:9.2.30}
设 $r>1$ 和 $\xi\in\mathrm{Pic}(C)$.
\begin{enumerate}[(a)]
\item 证明 $\cBun_{r,\xi}$ 是维数 $(r^2-1)(g-1)$ 的光滑代数叠.
\item 证明在 $E\in\cBun_{r,\xi}(\Bbbk)$ 处的切空间与映射 $\operatorname{Ext}^1(E,E)\to\mathrm H^1(C,\mathcal O_C)$ 的核等同; 此映射由迹映射 $E\otimes E^\vee\to\mathcal O_C$ 取 $\mathrm H^1$ 诱导.
\end{enumerate}
\end{exercise}
\begin{exercise}[困难]\label{src:9.2.31}
若 $(C,p)$ 是椭圆曲线, $\xi$ 是次数 $d$ 的线丛, 证明 $M_{r,\xi}^{\mathrm{ps}}\cong\mathbb P^{\gcd(r,d)-1}$. (这将习题~\bookref{9.1.47} 的分类推广到固定行列式的丛. 另见 \cite[定理 3]{Tu93}.)
\end{exercise}
\begin{remark}\label{src:9.2.32}
除了使用纤维图 \bookeqref{9.2.28} 定义固定行列式的向量丛叠, 也可使用剩余 gerbe: 若 $B\mathbb G_m\hookrightarrow\underline{\mathcal P\!ic}^d$ 是 $\xi$ 的剩余 gerbe 的包含映射, 则纤维积 $\mathcal X:=\cBun_{r,d}\times_{\underline{\mathcal P\!ic}^d}B\mathbb G_m$ 参数化秩为 $r$ 且存在同构 $\det E\xrightarrow{\sim}\xi$ 的向量丛 $E$. 稳定向量丛 $E\in\mathcal X(\Bbbk)$ 的自同构群为 $\mathbb G_m$, 稳定轨迹 $\mathcal X^{\mathrm s}:=\cBun_{r,d}^{\mathrm s}\times_{\underline{\mathcal P\!ic}^d}B\mathbb G_m$ 是如下空间上的 $\mathbb G_m$-gerbe:
\[
\mathcal X^{\mathrm s}\!\sslash\mathbb G_m\cong M_{r,\xi}^{\mathrm s}
\cong\cBun_{r,\xi}^{\mathrm s}\!\sslash\boldsymbol\mu_r.
\]
\end{remark}
